\documentclass[11pt]{article}

\usepackage[utf8]{inputenc}
\usepackage{amsmath, amssymb}
\usepackage{graphicx}
\usepackage{booktabs}
\usepackage[numbers]{natbib}
\usepackage{hyperref}

\title{Ink Density Saturation in Duan Inkstones:\\A Rendering Check}
\author{Harry Wang}
\date{August 2026}

\begin{document}

\maketitle

\begin{abstract}
A short sample document for testing \texttt{.tex} rendering, syntax
highlighting, and citation handling. The content mirrors the vault's
notes on grinding ink: density saturates as a function of grind time,
and the stone --- not the stick --- sets the ceiling.
\end{abstract}

\section{Introduction}

Grinding ink is a slow circular process: water is added \emph{a few
drops} at a time, and the stone decides how dark the ink gets. Prior
work established the density--time relationship empirically
\citep{tanaka2019}, while the classical treatment of stone provenance
goes back much further \citep{mizuta1968}. A modern survey is given by
\citet{chen2021}.

\section{Model}

Let $D(t)$ denote ink density after $t$ minutes of grinding. Density
follows a saturating exponential,
\begin{equation}
  \label{eq:density}
  D(t) = D_{\max}\left(1 - e^{-t/\tau}\right),
\end{equation}
where $D_{\max}$ is the stone-specific ceiling and $\tau$ the time
constant. Saturation is effectively reached near $t^* \approx 4$ min,
since from Equation~\eqref{eq:density},
\begin{align}
  \frac{D(t^*)}{D_{\max}} &= 1 - e^{-t^*/\tau} \\
                          &\approx 0.98 \quad \text{for } \tau \approx 1\ \text{min}.
\end{align}

Inline math also needs checking: the residual is $\varepsilon(t) =
D_{\max} e^{-t/\tau}$, and $\lim_{t \to \infty} \varepsilon(t) = 0$.

\section{Stones}

Table~\ref{tab:stones} lists the stones used in the vault's notes.

\begin{table}[h]
  \centering
  \begin{tabular}{llrl}
    \toprule
    Stone & Origin    & Year & Grind \\
    \midrule
    Duan  & Guangdong & 1782 & slow  \\
    She   & Anhui     & 1854 & fine  \\
    \bottomrule
  \end{tabular}
  \caption{Inkstones referenced in the grinding notes.}
  \label{tab:stones}
\end{table}

\section{Conclusion}

If this document compiles, renders with highlighting, and resolves all
three citations from \texttt{references.bib}, the check passes.

\bibliographystyle{plainnat}
\bibliography{references}

\end{document}
